\section{Tujuan Serangan (Hasil yang Ingin Dicapai)}
Selain berdasarkan informasi, serangan juga dapat dibedakan berdasarkan target akhirnya:
\begin{itemize}
    \item \textbf{Key-Recovery Attack}: Tujuan utamanya adalah memperoleh kunci kriptografi.
    \item \textbf{Plaintext-Recovery Attack}: Tujuan utamanya adalah memperoleh plainteks, meskipun kunci tidak ditemukan.
    \item \textbf{Distinguishing Attack}: Menentukan apakah suatu data adalah keluaran algoritma kriptografi atau data acak (\textit{random}).
    \item \textbf{Forgery Attack}: Membuat pesan palsu yang terlihat sah bagi sistem (sangat krusial pada MAC dan tanda tangan digital).
    \item \textbf{Replay Attack}: Mengirim kembali pesan sah yang pernah dicegat sebelumnya untuk memicu tindakan berulang tanpa perlu memecahkan enkripsi.
\end{itemize}

